\documentclass[10pt,a4paper]{amsart}
\usepackage[margin=19mm]{geometry}
\usepackage{amsmath,amssymb,mathtools}
\usepackage[scr=rsfs,cal=euler]{mathalfa}
\usepackage{xfrac}
\usepackage[unicode,hidelinks,bookmarks=false]{hyperref}
\newcommand{\ie}{i.e.\ }
\newcommand{\cf}{cf.\ }
\newcommand{\resp}{respectively\ }
\newcommand{\Subsection}{\subsection}
\providecommand{\nobreakdash}{\nobreak\hbox{-}}
\setlength{\emergencystretch}{2em}
\multlinegap0pt
\usepackage{amssymb}
\newtheorem{prop}{Proposition}
\newtheorem{theo}{Theorem}
\newcommand{\RR}{\mathbb{R}}
\def\beq{\begin{equation}}
\def\eeq{\end{equation}}
\title{Holomorphic Limits of K\"ahler Manifolds}
\author{Mu-Lin Li}
\date{Source: arXiv:2409.19957v4 (CC BY 4.0)}
\begin{document}
\maketitle

\noindent\emph{Source and licence.} Holomorphic Limits of K\"ahler Manifolds. Version arXiv:2409.19957v4 (CC BY 4.0), distributed by arXiv under the Creative Commons Attribution 4.0 International licence, http://creativecommons.org/licenses/by/4.0/, as recorded in the arXiv licence metadata for this exact version and shown on the abstract page https://arxiv.org/abs/2409.19957v4. Full licence notice: \texttt{licenses/arxiv-2409.19957-CC-BY-4.0.txt}.\par\smallskip

\noindent\textit{Editorial scope.} Counted unit: the article's prop good (source0.tex lines 634--641). The article's complete original proof of it is delivered with it (source0.tex lines 642--743, one \texttt{proof} environment of 102 lines). No mathematics is written, repaired or re-derived: the statement and the proof are the article's own lines, in the article's own notation, and no retained line is substituted or re-flowed.  Retained uncounted as the source-local context the proof needs: lines 592--597; lines 600--603; lines 624--631.  Nothing was omitted from any retained span, so the package carries no omission record.  The retained lines cite 3 works, whose entries are transcribed verbatim from the article's own bibliography source in the e-print and delivered with short editorial labels recorded in the item's reference mapping.  No retained line contains a figure environment, an image-inclusion command or drawing code, so no image is delivered and the package declares no asset dependency.  Editorial formatting only, in addition: an AMS \texttt{amsart} preamble that restates the article's own declarations and macros used by the retained lines, namely the environments \texttt{prop}, \texttt{theo} on separate counters, together with the standard packages those lines need.  The retained environments print in this document as Proposition 1, Theo 1 in order of appearance, whereas the article prints the same items with the numbering of its own full document.  The complete native source of the article as distributed in the arXiv e-print for this exact version is bundled unchanged as \texttt{sources/166020/source0.tex}.
\begin{prop}[{\cite[Proposition 3.16]{BGL}}]\label{equal}
If $\theta$ is a smooth $d$-closed $(1,1)$-form and $\{\theta\}$ is pseudo-effective, then
\beq
\underline{\mathrm{vol}}_n(\{\theta\})=\overline{\mathrm{vol}}_n(\{\theta\}).\nonumber
\eeq
\end{prop}

\begin{theo}[{\cite[Theorem 3.20]{BGL}}]\label{vanishing-vol}
The volume function $\mathrm{vol}_n:\allowbreak H^{1,1}_{BC}(X,\mathbb{R}) \to \mathbb{R}_{\ge0}$ is continuous. In
particular, it vanishes outside of the big cone.
\end{theo}

\begin{theo}[{\cite[Proposition 3.7]{Dem92},\cite[Theorem 2.4]{Bou}}]\label{regularization}
Let $T=\theta+dd^c\varphi$ be a $d$-closed positive $(1,1)$-current on a compact Hermitian manifold $(X,\omega_X)$, where $\theta$ is a smooth real $(1,1)$-form and $\varphi$ is quasi-psh. Assume that $T\ge\gamma$ for a smooth real $(1,1)$-form $\gamma$. Then there exists a sequence of functions $\varphi_k$ with analytic singularities converging to $\varphi$ such that, with $T_k:=\theta+dd^c\varphi_k$, the following hold:
\begin{itemize}
\item $T_k\to T$ weakly;
\item $T_k\ge\gamma-\epsilon_k\omega_X$, where $\epsilon_k>0$ and $\epsilon_k\to0$;
\item $T_{k,ac}\to T_{ac}$ almost everywhere on $X$
\end{itemize}
\end{theo}

\begin{prop}\label{good}
Let $X$ be an $n$-dimensional compact complex manifold with a Hermitian metric
$\omega_X$ satisfying $\overline{\mathrm{vol}}_n(\omega_X)<\infty$. Let
$\theta$ be a $d$-closed smooth real $(1,1)$-form and assume that $\{\theta\}\in H^{1,1}_{BC}(X,\RR)$ is pseudo-effective. Then
\[
 \mathrm{vol}_n(\{\theta\})=\widetilde{\mathrm{vol}}_n(\{\theta\}).
\]
\end{prop}

\begin{proof}
We prove the two inequalities separately. The finiteness assumption ensures that all volumes below are finite; this also follows from the estimates established in the first part of the argument.

Let $T=\theta+dd^c\varphi\ge0$ be an arbitrary $d$-closed positive current in
$\{\theta\}$. Apply Theorem~\ref{regularization} and write
\[
 T_k=\theta+dd^c\varphi_k,\qquad
 T_k\ge-\varepsilon_k\omega_X,\qquad
 \varepsilon_k\downarrow0,
\]
with $T_{k,ac}\to T_{ac}$ almost everywhere. Put
$S_k=T_k+\varepsilon_k\omega_X\ge0$. Resolving the analytic singularities of
$T_k$, we obtain a modification $\rho_k:Y_k\to X$ such that
\[
 \rho_k^*T_k=\beta_k+[E_k],
\]
where $\beta_k$ is smooth and $E_k$ is an effective divisor. Hence
\[
 (\rho_k^*S_k)_{ac}=\beta_k+\varepsilon_k\rho_k^*\omega_X\ge0
\]
and, since the exceptional and polar sets have Lebesgue measure zero,
\begin{equation}\label{eq:regularized-ac-mass}
 \int_X(T_{k,ac}+\varepsilon_k\omega_X)^n
 =\int_{Y_k}(\beta_k+\varepsilon_k\rho_k^*\omega_X)^n.
\end{equation}
The invariance of the upper volume under modifications and addition of divisors, together with the definition for smooth positive forms, gives
\begin{equation}\label{eq:upper-controls-ac}
 \overline{\mathrm{vol}}_n(\{\theta+\varepsilon_k\omega_X\})
 \ge \int_X(T_{k,ac}+\varepsilon_k\omega_X)^n;
\end{equation}
see \cite[Lemma 3.5 and Definition 3.8]{BGL}. Since the forms on the right-hand
side are semipositive and converge almost everywhere to $T_{ac}$, Fatou's lemma
yields
\[
 \int_XT_{ac}^n
 \le\liminf_{k\to\infty}
 \int_X(T_{k,ac}+\varepsilon_k\omega_X)^n.
\]
By Definition 3.13 of \cite{BGL} and Proposition~\ref{equal},
\[
 \lim_{k\to\infty}
 \overline{\mathrm{vol}}_n(\{\theta+\varepsilon_k\omega_X\})
 =\mathrm{vol}_n(\{\theta\}).
\]
Combining this with \eqref{eq:upper-controls-ac}, we obtain
\[
 \int_XT_{ac}^n\le\mathrm{vol}_n(\{\theta\}).
\]
Taking the supremum over all $T\ge0$ in $\{\theta\}$ gives
\begin{equation}\label{eq:tilde-le-volume}
 \widetilde{\mathrm{vol}}_n(\{\theta\})
 \le\mathrm{vol}_n(\{\theta\}).
\end{equation}
If the right-hand side vanishes, the desired equality follows immediately.

Assume now that $\mathrm{vol}_n(\{\theta\})>0$. By
Theorem~\ref{vanishing-vol}, the class $\{\theta\}$ is big. Fix $\delta>0$.
Proposition 3.12 of \cite{BGL} gives a modification $\rho:Y\to X$ and a
decomposition
\begin{equation}\label{eq:BGL-decomposition}
 \{\rho^*\theta\}=\{\omega_Y\}+\{E\},
\end{equation}
where $\omega_Y>0$ is smooth, $E$ is an effective real divisor, and
\begin{equation}\label{eq:BGL-near-volume}
 \underline{\mathrm{vol}}_n(\omega_Y)
 >\mathrm{vol}_n(\{\theta\})-\delta.
\end{equation}
We claim that $\omega_Y$ is $d$-closed. Indeed, equality
\eqref{eq:BGL-decomposition} means that
$\rho^*\theta-\omega_Y-[E]=dd^cu$ for a distribution $u$; applying $d$ and
using $d\theta=d[E]=0$ gives $d\omega_Y=0$. Thus $\omega_Y$ is a K\"ahler
form. Proposition 3.16 of \cite{BGL} therefore gives
\begin{equation}\label{eq:omegaY-volume}
 \underline{\mathrm{vol}}_n(\omega_Y)=\int_Y\omega_Y^n.
\end{equation}

Set
\[
 R:=\rho_*(\omega_Y+[E]).
\]
Then $R$ is a $d$-closed positive $(1,1)$-current on $X$. Pushing forward
\eqref{eq:BGL-decomposition} and using $\rho_*\rho^*\theta=\theta$ for a
modification shows that $\{R\}=\{\theta\}$. Since
$(\omega_Y+[E])_{ac}=\omega_Y$, Boucksom's
push-forward formula for absolutely continuous parts
\cite[Proposition 2.2]{Bou} gives
\[
 R_{ac}=\rho_*\omega_Y\qquad\text{almost everywhere}.
\]
A modification is biholomorphic outside analytic subsets of Lebesgue measure
zero. Consequently
\begin{equation}\label{eq:push-ac-mass}
 \int_XR_{ac}^n=\int_Y\omega_Y^n.
\end{equation}
Equations \eqref{eq:BGL-near-volume}--\eqref{eq:push-ac-mass} imply
\[
 \widetilde{\mathrm{vol}}_n(\{\theta\})
 \ge\int_XR_{ac}^n
 >\mathrm{vol}_n(\{\theta\})-\delta.
\]
Letting $\delta\downarrow0$ and combining the resulting inequality with \eqref{eq:tilde-le-volume} proves the proposition.
\end{proof}

\begin{thebibliography}{99}
\bibitem[BGL]{BGL} S. Boucksom, V. Guedj, and C. H. Lu, \textit{Volumes of Bott--Chern classes}, Peking Math. J. (2025), published online 16 June 2025, doi:10.1007/s42543-025-00105-2.
\bibitem[Bou]{Bou} S. Boucksom, \textit{On the volume of a line bundle}, Internat. J. Math. \textbf{13} (2002), no.~10, 1043--1063.
\bibitem[Dem92]{Dem92} J.-P. Demailly, \textit{Regularization of closed positive currents and intersection theory}, J. Algebraic Geom. \textbf{1} (1992), 361--409.
\end{thebibliography}
\end{document}
